\documentclass{article}

% Language setting
% Replace `english' with e.g. `spanish' to change the document language
\usepackage[czech]{babel}
\usepackage{amsthm,thmtools,xcolor,amsmath,amssymb, mathtools}

% Set page size and margins
% Replace `letterpaper' with `a4paper' for UK/EU standard size
\usepackage[a4paper,top=2cm,bottom=2cm,left=3cm,right=3cm,marginparwidth=1.75cm]{geometry}

% Useful packages
\usepackage{amsmath}
\usepackage{enumerate}
\usepackage{graphicx}
\usepackage[colorlinks=true, allcolors=blue]{hyperref}
\graphicspath{images/}

%\usepackage{tikzit}
%\input{default.tikzstyles}

\title{DÚ Lineární algebra -- Sada 5}
\author{Jan Romanovský}

\begin{document}
\maketitle

\textbf{(5.2)} Podprostor příslušný vlastnímu číslu 2 je geometricky rovina, tedy prostor dimenze 2, vyberu si z ní tedy dva LN vektory, např. $\left(1, 1, 0\right)^T$, $\left(0, 1, 1\right)^T$. Ty pak spolu s vektorem $\left(1, 1, 1\right)^T$ tvoří LN bázi a matice operátoru v této bázi bude:
$$A_B = \begin{pmatrix}
   2 & 0 & 0 \\ 0 & 2 & 0 \\ 0 & 0 & 3
\end{pmatrix}$$
Matice přechodu vidíme přímo:
$$[\text{id}]_K^B = \begin{pmatrix}
    1 & 0 & 1 \\ 1 & 1 & 1 \\ 0 & 1 & 1
\end{pmatrix},\, [\text{id}]_B^K = ([\text{id}]_K^B)^{-1} = \begin{pmatrix}
   0 & 1 & -1 \\ -1 & 1 & 0 \\ 1 & -1 & 1
\end{pmatrix}$$
Matice operátoru v kanonické bázi je pak:
$$A = [\text{id}]_K^BA_B[\text{id}]_B^K = \begin{pmatrix}
    3 & -1 & 1 \\ 1 & 1 & 1 \\ 1 & -1 & 3
\end{pmatrix}$$
\end{document}

